\section{Part F. Topological Sort}
\begin{frame}{Topological Order}
\[
(u,v)\in E\quad\Longrightarrow\quad pos[u]<pos[v].
\]
\begin{block}{필요충분조건}directed graph가 topological order를 갖는다 \(\Longleftrightarrow\) graph가 DAG다.\end{block}
\begin{itemize}\item valid order는 여러 개일 수 있다.\item tie/container policy에 따라 concrete order가 달라진다.\end{itemize}
\end{frame}
\begin{frame}{Kahn Algorithm}
\scriptsize\begin{algorithmic}[1]
\Procedure{TopoKahn}{$G$}\State 모든 \(indegree[v]\) 계산; \(Q\gets\) indegree 0 vertices; \(order\gets[]\)
\While{$Q\ne\emptyset$}\State \(u\gets\Call{Remove}{Q}\); append \(u\)
\ForAll{$v\in Adj[u]$}\State \(indegree[v]\gets indegree[v]-1\)
\If{$indegree[v]=0$}\State \Call{Insert}{$Q,v$}\EndIf\EndFor\EndWhile
\If{$|order|\ne|V|$}\State\Return CYCLE\EndIf\State\Return order\EndProcedure
\end{algorithmic}
\end{frame}
\begin{frame}{Kahn의 선택 정책}
\begin{columns}[T]\column{.32\textwidth}\begin{block}{Queue}발견된 순서\\\(\Theta(V+E)\)\end{block}
\column{.32\textwidth}\begin{block}{Stack}최근 zero 먼저\\\(\Theta(V+E)\)\end{block}
\column{.32\textwidth}\begin{block}{Min-heap}lexicographically smallest\\\(O(E+V\log V)\)\end{block}\end{columns}
\gtakeaway{“random node”가 아니라 현재 zero-indegree container에서 정책에 따라 하나를 선택한다.}
\end{frame}
\begin{frame}{Kahn Trace: Min-Heap}
\small edges: \(0\to1,0\to3,1\to2,1\to4,2\to5,3\to5,4\to5\).
\centering\begin{tikzpicture}
\node[pq]at(0,1){Zero heap: \only<1>{[0]}\only<2>{[1,3]}\only<3>{[2,3,4]}\only<4>{[3,4]}\only<5>{[4]}\only<6>{[5]\(\to\)[]}};
\node[callout]at(0,0){Select: \only<1>{--}\only<2>{0}\only<3>{1}\only<4>{2}\only<5>{3}\only<6>{4 then 5}};
\node[callout]at(0,-1){Output: \only<1>{[]}\only<2>{[0]}\only<3>{[0,1]}\only<4>{[0,1,2]}\only<5>{[0,1,2,3]}\only<6>{[0,1,2,3,4,5]}};
\end{tikzpicture}
\end{frame}
\begin{frame}{Kahn Cycle Detection}
\centering\begin{tikzpicture}[x=1.5cm]
\node[gvertex](a)at(0,0){A};\node[gvertex](b)at(2,0){B};\node[gvertex](c)at(4,0){C};
\draw[garrow](a)--(b);\draw[garrow](b)--(c);\draw[gback](c)to[bend left=35](a);
\end{tikzpicture}
\begin{alertblock}{판정}zero-indegree container가 비고 \(processedCount<|V|\)이면 directed cycle. partial output을 full order로 반환하지 않는다.\end{alertblock}
\end{frame}
\begin{frame}{DFS Topological Sort}
\scriptsize\begin{algorithmic}[1]
\Function{TopoVisit}{$u$}\State \(color[u]\gets G\)
\ForAll{$v\in Adj[u]$}\If{$color[v]=G$}\State\Return CYCLE\EndIf
\If{$color[v]=W$ \textbf{ and } \Call{TopoVisit}{$v$}=\text{CYCLE}}\State\Return CYCLE\EndIf\EndFor
\State \(color[u]\gets B\); append \(u\) to \(order\); \Return OK\EndFunction
\Statex outer loop 후 \(\Call{Reverse}{order}\)
\end{algorithmic}
\end{frame}
\begin{frame}{Finish Order의 Reverse}
\centering\begin{tikzpicture}
\node[pq]at(0,.8){Finish append: \only<1>{5}\only<2>{5,2}\only<3>{5,2,4,1}\only<4>{5,2,4,1,3,0}};
\node[callout]at(0,-.3){Reverse: \only<1-3>{아직 진행 중}\only<4>{0,3,1,4,2,5}};
\end{tikzpicture}
\begin{block}{이유}edge \(u\to v\)에서 cycle이 없다면 \(v\)의 DFS가 끝난 뒤 \(u\)가 끝나므로 \(f[u]>f[v]\). decreasing finish time은 edge 방향을 존중한다.\end{block}
\end{frame}
\begin{frame}{Kahn vs DFS Topological Sort}
\scriptsize\centering\begin{tabular}{lll}\toprule&Kahn&DFS\\\midrule
state&indegree&color, finish\\cycle&processed \(<V\)&GRAY back edge\\order control&queue/stack/min-heap&outer/adj order\\storage&container&recursion stack\\
time: queue/stack&\(\Theta(V+E)\)&\(\Theta(V+E)\)\\
time: min-heap&\(O(E+V\log V)\)&N/A\\\bottomrule\end{tabular}
\begin{block}{Heap cost}각 vertex는 zero-indegree heap에 최대 한 번 들어가고 한 번 extract된다.\end{block}
\end{frame}
\begin{frame}{활용과 한계}
\begin{itemize}\item prerequisite, build/package dependency\item task dependency, DAG shortest path\item Lecture 5 DP evaluation order\end{itemize}
\begin{alertblock}{주의}topological order만으로 resource와 duration을 고려한 최소 parallel schedule이 자동으로 나오지는 않는다.\end{alertblock}
\end{frame}
\begin{frame}{Part F Checkpoint}
\begin{enumerate}\item topo order의 존재 조건은?\item lexicographically smallest order에 쓸 container는?\item 두 방법의 cycle signal은?\item min-heap Kahn의 list-graph 시간은?\end{enumerate}
\begin{exampleblock}{답}DAG; min-heap; Kahn은 processed count, DFS는 GRAY back edge; \(O(E+V\log V)\).\end{exampleblock}
\end{frame}
